% =============================================================================
% CHAPTER 18: TEXT-TO-SQL DRIFT DETECTION FRAMEWORK
% =============================================================================

\chapter{Text-to-SQL Drift Detection Framework}
\label{chap:tts-drift}

\section{Purpose and Scope}

This chapter defines a comprehensive drift detection framework for Text-to-SQL
pipelines. Drift occurs when the relationship between any of the following
artifacts diverges from the expected baseline:

\begin{enumerate}
  \item \textbf{HANA Table Schemas} --- The source of truth for database structure
  \item \textbf{Agent/User Prompts} --- Natural language queries from users or agents
  \item \textbf{Generated SQL} --- The Text-to-SQL model output
\end{enumerate}

The framework integrates with:
\begin{itemize}
  \item Simula training data generation (Chapter~\ref{chap:generator})
  \item Spec-drift audit infrastructure (Chapter~4a of clinerules-agents)
  \item Singapore AI Safety metrics framework (Chapter~5 of singapore-ai-safety)
\end{itemize}

\section{Drift Taxonomy}

The Text-to-SQL drift taxonomy distinguishes three primary drift sources:

\begin{enumerate}
  \item \textbf{Schema Drift} --- Changes to HANA database structure not reflected in training
  \item \textbf{User Prompt Drift} --- Evolution of human user query patterns and vocabulary
  \item \textbf{Agent Prompt Drift} --- Changes in AI agent query generation behavior
\end{enumerate}

\begin{table}[htbp]
\centering
\footnotesize
\caption{Text-to-SQL Drift Types by Source}
\label{tab:tts-drift-types}
\begin{tabularx}{\textwidth}{@{}l l l X l@{}}
\toprule
\textbf{Code} & \textbf{Type} & \textbf{Source} & \textbf{Description} & \textbf{Severity} \\
\midrule
TTS-DRIFT-001 & Schema Drift & Schema & HANA schema changed without training update & CRITICAL \\
TTS-DRIFT-002 & Coverage Drift & Schema & Training no longer covers schema adequately & HIGH \\
TTS-DRIFT-003 & Semantic Drift & User/Agent & Prompt intent no longer aligns with SQL & CRITICAL \\
TTS-DRIFT-004 & Terminology Drift & User & User vocabulary diverged from schema terms & MEDIUM \\
TTS-DRIFT-005 & Complexity Drift & User/Agent & Query complexity differs from training & MEDIUM \\
TTS-DRIFT-006 & Execution Drift & All & SQL execution success rate degraded & CRITICAL \\
TTS-DRIFT-007 & Agent Behavior Drift & Agent & Agent query patterns changed significantly & HIGH \\
TTS-DRIFT-008 & User Behavior Drift & User & User query patterns changed significantly & MEDIUM \\
\bottomrule
\end{tabularx}
\end{table}

\subsection{Schema Drift Sources (TTS-DRIFT-001, TTS-DRIFT-002)}

\textbf{Definition:} Schema drift occurs when the HANA Cloud database structure
changes and the Text-to-SQL training data becomes stale or inconsistent.

\begin{table}[htbp]
\centering
\footnotesize
\caption{Schema Drift Trigger Events}
\label{tab:schema-drift-triggers}
\begin{tabularx}{\textwidth}{@{}l l l X@{}}
\toprule
\textbf{Event} & \textbf{Weight} & \textbf{Impact} & \textbf{Mitigation} \\
\midrule
Column dropped & 1.0 & SQL references fail & Immediate retraining \\
Column type changed & 0.5 & Type errors in queries & Schema re-extraction \\
Column renamed & 0.3 & Minor query failures & Vocabulary update \\
Table added & 0.2 & Missing coverage & Incremental training \\
FK relationship changed & 0.4 & Invalid JOINs & Schema refresh \\
\bottomrule
\end{tabularx}
\end{table}

\textbf{Rating Criteria for Schema Drift Metrics:}

\begin{table}[htbp]
\centering
\footnotesize
\caption{Schema Drift Metric Quality Assessment}
\label{tab:schema-metric-rating}
\begin{tabularx}{\textwidth}{@{}l c c c c X@{}}
\toprule
\textbf{Metric} & \textbf{Accuracy} & \textbf{Coverage} & \textbf{Speed} & \textbf{Overall} & \textbf{Notes} \\
\midrule
TTS-M01 (SCR) & 95\% & High & Fast & \textbf{A} & Direct column comparison \\
TTS-M02 (SSS) & 90\% & High & Fast & \textbf{A} & Time-weighted staleness \\
TTS-M03 (CTMR) & 85\% & Medium & Medium & \textbf{B+} & Requires SQL parsing \\
TTS-M04 (FKCR) & 90\% & High & Medium & \textbf{A-} & FK metadata required \\
\bottomrule
\end{tabularx}
\end{table}

\subsection{User Prompt Drift Sources (TTS-DRIFT-004, TTS-DRIFT-008)}

\textbf{Definition:} User prompt drift occurs when human users change their
query patterns, vocabulary, or complexity expectations over time.

\begin{table}[htbp]
\centering
\footnotesize
\caption{User Prompt Drift Indicators}
\label{tab:user-drift-indicators}
\begin{tabularx}{\textwidth}{@{}l X l l@{}}
\toprule
\textbf{Indicator} & \textbf{Description} & \textbf{Metric} & \textbf{Threshold} \\
\midrule
Vocabulary shift & New terms not in training data & TTS-M07 & JSD $> 0.15$ \\
Query complexity increase & Users asking harder questions & TTS-M11 & KL $> 0.10$ \\
Ambiguity increase & More vague or underspecified prompts & TTS-M08 & ARR $< 0.75$ \\
Domain expansion & Users querying new table areas & TTS-M01 & SCR $< 0.85$ \\
Intent mismatch & Users expect different results & TTS-M06 & IPR $< 0.90$ \\
\bottomrule
\end{tabularx}
\end{table}

\textbf{Rating Criteria for User Prompt Drift Metrics:}

\begin{table}[htbp]
\centering
\footnotesize
\caption{User Prompt Drift Metric Quality Assessment}
\label{tab:user-metric-rating}
\begin{tabularx}{\textwidth}{@{}l c c c c X@{}}
\toprule
\textbf{Metric} & \textbf{Accuracy} & \textbf{Coverage} & \textbf{Speed} & \textbf{Overall} & \textbf{Notes} \\
\midrule
TTS-M05 (SAS) & 80\% & High & Slow & \textbf{B+} & Embedding-dependent \\
TTS-M06 (IPR) & 85\% & High & Slow & \textbf{B+} & Requires critic model \\
TTS-M07 (TDS) & 90\% & Medium & Fast & \textbf{A-} & JSD computation \\
TTS-M08 (ARR) & 75\% & Medium & Slow & \textbf{B} & Ambiguity detection hard \\
\bottomrule
\end{tabularx}
\end{table}

\subsection{Agent Prompt Drift Sources (TTS-DRIFT-003, TTS-DRIFT-007)}

\textbf{Definition:} Agent prompt drift occurs when AI agents (Cline, copilots,
or automated systems) change their query generation behavior.

\begin{table}[htbp]
\centering
\footnotesize
\caption{Agent Prompt Drift Indicators}
\label{tab:agent-drift-indicators}
\begin{tabularx}{\textwidth}{@{}l X l l@{}}
\toprule
\textbf{Indicator} & \textbf{Description} & \textbf{Metric} & \textbf{Threshold} \\
\midrule
Template evolution & Agent prompts follow new patterns & TTS-M05 & SAS $< 0.80$ \\
Instruction compliance & Agent adheres to system prompts & TTS-M06 & IPR $< 0.90$ \\
Query structure shift & SQL generation patterns changed & TTS-M11 & CDD $> 0.10$ \\
Tool calling patterns & Different tool invocation sequences & TTS-M09 & ESR $< 0.95$ \\
Context window usage & Agent uses context differently & TTS-M10 & RFS $< 0.85$ \\
\bottomrule
\end{tabularx}
\end{table}

\textbf{Rating Criteria for Agent Prompt Drift Metrics:}

\begin{table}[htbp]
\centering
\footnotesize
\caption{Agent Prompt Drift Metric Quality Assessment}
\label{tab:agent-metric-rating}
\begin{tabularx}{\textwidth}{@{}l c c c c X@{}}
\toprule
\textbf{Metric} & \textbf{Accuracy} & \textbf{Coverage} & \textbf{Speed} & \textbf{Overall} & \textbf{Notes} \\
\midrule
TTS-M05 (SAS) & 85\% & High & Slow & \textbf{A-} & Effective for agent drift \\
TTS-M06 (IPR) & 90\% & High & Slow & \textbf{A} & Critical for agents \\
TTS-M09 (ESR) & 95\% & High & Fast & \textbf{A+} & Direct execution check \\
TTS-M10 (RFS) & 85\% & Medium & Slow & \textbf{B+} & Result comparison \\
TTS-M11 (CDD) & 80\% & Medium & Fast & \textbf{B+} & Distribution comparison \\
\bottomrule
\end{tabularx}
\end{table}

\subsection{Distinguishing User vs Agent Prompts}

The framework must distinguish between human user prompts and AI agent prompts
for accurate drift attribution:

\begin{table}[htbp]
\centering
\footnotesize
\caption{User vs Agent Prompt Classification}
\label{tab:user-agent-classification}
\begin{tabularx}{\textwidth}{@{}l X X@{}}
\toprule
\textbf{Characteristic} & \textbf{User Prompts} & \textbf{Agent Prompts} \\
\midrule
Variability & High (natural language variety) & Low (templated patterns) \\
Vocabulary & Domain-specific, informal & Precise, instruction-following \\
Complexity & Variable (beginner to expert) & Consistent (system-configured) \\
Ambiguity & Common (underspecified) & Rare (explicit requirements) \\
Context usage & Limited (single query) & Rich (multi-turn, tools) \\
Error patterns & Typos, vague terms & Tool calling errors, format issues \\
\bottomrule
\end{tabularx}
\end{table}

\textbf{Classification Method:}

\begin{lstlisting}[language=jsonx,caption={User vs Agent prompt classification}]
def classify_prompt_source(prompt: str, metadata: dict) -> str:
    """
    Classify prompt as user or agent based on characteristics.
    
    Returns: "user", "agent", or "unknown"
    """
    # Check metadata first
    if "agent_id" in metadata or "tool_call_id" in metadata:
        return "agent"
    
    if "user_session_id" in metadata:
        # Additional heuristics for user prompts
        features = {
            "has_typos": detect_typos(prompt),
            "uses_informal_language": detect_informal(prompt),
            "high_variability": compute_perplexity(prompt) > 10,
            "explicit_instructions": has_system_instructions(prompt),
        }
        
        user_score = (features["has_typos"] + 
                      features["uses_informal_language"] +
                      features["high_variability"])
        agent_score = features["explicit_instructions"]
        
        return "user" if user_score > agent_score else "agent"
    
    return "unknown"
\end{lstlisting}

\subsection{Combined Drift Assessment Matrix}

\begin{table}[htbp]
\centering
\footnotesize
\caption{Comprehensive Drift Assessment Matrix}
\label{tab:drift-assessment-matrix}
\begin{tabularx}{\textwidth}{@{}l c c c c@{}}
\toprule
\textbf{Metric} & \textbf{Schema} & \textbf{User Prompt} & \textbf{Agent Prompt} & \textbf{Weight} \\
\midrule
TTS-M01 (SCR) & \checkmark & -- & -- & 15\% \\
TTS-M02 (SSS) & \checkmark & -- & -- & 10\% \\
TTS-M03 (CTMR) & \checkmark & -- & -- & 5\% \\
TTS-M04 (FKCR) & \checkmark & -- & -- & 10\% \\
TTS-M05 (SAS) & -- & \checkmark & \checkmark & 15\% \\
TTS-M06 (IPR) & -- & \checkmark & \checkmark & 10\% \\
TTS-M07 (TDS) & -- & \checkmark & -- & 3\% \\
TTS-M08 (ARR) & -- & \checkmark & -- & 5\% \\
TTS-M09 (ESR) & \checkmark & \checkmark & \checkmark & 15\% \\
TTS-M10 (RFS) & -- & \checkmark & \checkmark & 10\% \\
TTS-M11 (CDD) & -- & \checkmark & \checkmark & 2\% \\
TTS-M12 (TCD) & \checkmark & -- & -- & -- \\
\bottomrule
\end{tabularx}
\end{table}

\textbf{Reading the matrix:}
\begin{itemize}
  \item \checkmark = Metric is relevant for this drift source
  \item Weight = Contribution to composite TTS-EVAL score
  \item Metrics can be relevant to multiple sources
\end{itemize}

\section{Mandatory Drift Metrics}

\subsection{Metric Overview}

\begin{table}[htbp]
\centering
\footnotesize
\caption{Text-to-SQL Drift Metrics (12 Metrics)}
\label{tab:tts-metrics}
\begin{tabularx}{\textwidth}{@{}l l X l@{}}
\toprule
\textbf{ID} & \textbf{Name} & \textbf{Definition} & \textbf{Drift Link} \\
\midrule
\metricid{TTS-M01} & Schema Coverage Rate (SCR) & Fraction of HANA columns referenced in training & TTS-DRIFT-002 \\
\metricid{TTS-M02} & Schema Staleness Score (SSS) & Age-weighted schema changes not in training & TTS-DRIFT-001 \\
\metricid{TTS-M03} & Column Type Mismatch Rate (CTMR) & Fraction of SQL with type-incompatible ops & TTS-DRIFT-001 \\
\metricid{TTS-M04} & Foreign Key Consistency Rate (FKCR) & Fraction of JOINs following valid FKs & TTS-DRIFT-001 \\
\metricid{TTS-M05} & Semantic Alignment Score (SAS) & Cosine similarity prompt-to-SQL embeddings & TTS-DRIFT-003 \\
\metricid{TTS-M06} & Intent Preservation Rate (IPR) & Fraction of prompts correctly answered & TTS-DRIFT-003 \\
\metricid{TTS-M07} & Terminology Drift Score (TDS) & Vocabulary divergence from schema terms & TTS-DRIFT-004 \\
\metricid{TTS-M08} & Ambiguity Resolution Rate (ARR) & Fraction of ambiguous prompts disambiguated & TTS-DRIFT-003 \\
\metricid{TTS-M09} & Execution Success Rate (ESR) & Fraction of SQL executing without error & TTS-DRIFT-006 \\
\metricid{TTS-M10} & Result Fidelity Score (RFS) & Semantic similarity expected vs actual results & TTS-DRIFT-006 \\
\metricid{TTS-M11} & Complexity Distribution Drift (CDD) & KL divergence training vs production complexity & TTS-DRIFT-005 \\
\metricid{TTS-M12} & Taxonomy Coverage Drift (TCD) & Change in taxonomy node coverage over time & TTS-DRIFT-002 \\
\bottomrule
\end{tabularx}
\end{table}

\subsection{Schema Drift Metrics (TTS-M01 to TTS-M04)}

\subsubsection{TTS-M01: Schema Coverage Rate (SCR)}

\begin{equation}
\boxed{
SCR = \frac{|C_{\text{referenced}}|}{|C_{\text{schema}}|}
}
\end{equation}

\begin{description}
  \item[Variables] $C_{\text{referenced}}$ = columns referenced in training examples;
        $C_{\text{schema}}$ = columns in HANA schema
  \item[Direction] Higher is better
  \item[Threshold] $SCR \geq 0.85$ required
  \item[Drift Link] \driftid{TTS-DRIFT-002}
\end{description}

\subsubsection{TTS-M02: Schema Staleness Score (SSS)}

\begin{equation}
\boxed{
SSS = \sum_{i \in \Delta_{\text{schema}}} w_i \cdot \text{age}(i)
}
\end{equation}

\begin{description}
  \item[Variables] $\Delta_{\text{schema}}$ = set of schema changes not in training;
        $w_i$ = importance weight (1.0 column drop, 0.5 type change, 0.3 rename);
        $\text{age}(i)$ = days since change
  \item[Direction] Lower is better
  \item[Threshold] $SSS \leq 7.0$ required (approximately 1 week tolerance)
  \item[Drift Link] \driftid{TTS-DRIFT-001}
\end{description}

\subsubsection{TTS-M03: Column Type Mismatch Rate (CTMR)}

\begin{equation}
\boxed{
CTMR = \frac{N_{\text{type-incompatible operations}}}{N_{\text{total column references}}}
}
\end{equation}

\begin{description}
  \item[Direction] Lower is better
  \item[Threshold] $CTMR \leq 0.02$ required
  \item[Drift Link] \driftid{TTS-DRIFT-001}
\end{description}

\subsubsection{TTS-M04: Foreign Key Consistency Rate (FKCR)}

\begin{equation}
\boxed{
FKCR = \frac{N_{\text{valid FK joins}}}{N_{\text{total joins}}}
}
\end{equation}

\begin{description}
  \item[Direction] Higher is better
  \item[Threshold] $FKCR \geq 0.95$ required
  \item[Drift Link] \driftid{TTS-DRIFT-001}
\end{description}

\subsection{Semantic Drift Metrics (TTS-M05 to TTS-M08)}

\subsubsection{TTS-M05: Semantic Alignment Score (SAS)}

\begin{equation}
\boxed{
SAS = \frac{1}{N} \sum_{j=1}^{N} \cos(E_{\text{prompt}}^j, E_{\text{sql}}^j)
}
\end{equation}

\begin{description}
  \item[Variables] $E_{\text{prompt}}$ = embedding of natural language prompt;
        $E_{\text{sql}}$ = embedding of SQL intent (parsed query structure);
        $\cos()$ = cosine similarity
  \item[Direction] Higher is better
  \item[Threshold] $SAS \geq 0.80$ required
  \item[Drift Link] \driftid{TTS-DRIFT-003}
\end{description}

\subsubsection{TTS-M06: Intent Preservation Rate (IPR)}

\begin{equation}
\boxed{
IPR = \frac{N_{\text{correctly answered prompts}}}{N_{\text{total prompts evaluated}}}
}
\end{equation}

\begin{description}
  \item[Direction] Higher is better
  \item[Threshold] $IPR \geq 0.90$ required
  \item[Drift Link] \driftid{TTS-DRIFT-003}
\end{description}

\subsubsection{TTS-M07: Terminology Drift Score (TDS)}

\begin{equation}
\boxed{
TDS = JSD(P_{\text{prompt\_vocab}} \| P_{\text{schema\_vocab}})
}
\end{equation}

\begin{description}
  \item[Variables] $P_{\text{prompt\_vocab}}$ = term probability distribution in prompts;
        $P_{\text{schema\_vocab}}$ = term distribution in schema (columns, tables);
        $JSD$ = Jensen-Shannon Divergence
  \item[Direction] Lower is better
  \item[Threshold] $TDS \leq 0.15$ required
  \item[Drift Link] \driftid{TTS-DRIFT-004}
\end{description}

\subsubsection{TTS-M08: Ambiguity Resolution Rate (ARR)}

\begin{equation}
\boxed{
ARR = \frac{N_{\text{ambiguous prompts correctly resolved}}}{N_{\text{ambiguous prompts total}}}
}
\end{equation}

\begin{description}
  \item[Direction] Higher is better
  \item[Threshold] $ARR \geq 0.75$ required
  \item[Drift Link] \driftid{TTS-DRIFT-003}
\end{description}

\subsection{Generation Quality Drift Metrics (TTS-M09 to TTS-M12)}

\subsubsection{TTS-M09: Execution Success Rate (ESR)}

\begin{equation}
\boxed{
ESR = \frac{N_{\text{SQL executed successfully}}}{N_{\text{SQL generated}}}
}
\end{equation}

\begin{description}
  \item[Direction] Higher is better
  \item[Threshold] $ESR \geq 0.95$ required
  \item[Drift Link] \driftid{TTS-DRIFT-006}
\end{description}

\subsubsection{TTS-M10: Result Fidelity Score (RFS)}

\begin{equation}
\boxed{
RFS = \frac{1}{N} \sum_{j=1}^{N} \text{sim}(R_{\text{expected}}^j, R_{\text{actual}}^j)
}
\end{equation}

\begin{description}
  \item[Variables] $R_{\text{expected}}$ = expected query result (ground truth);
        $R_{\text{actual}}$ = actual query result; $\text{sim}()$ = semantic similarity
  \item[Direction] Higher is better
  \item[Threshold] $RFS \geq 0.85$ required
  \item[Drift Link] \driftid{TTS-DRIFT-006}
\end{description}

\subsubsection{TTS-M11: Complexity Distribution Drift (CDD)}

\begin{equation}
\boxed{
CDD = KL(P_{\text{train\_complexity}} \| P_{\text{prod\_complexity}})
}
\end{equation}

\begin{description}
  \item[Variables] Using Elo-calibrated complexity scores from Algorithm~3
        (Chapter~\ref{chap:evaluation})
  \item[Direction] Lower is better
  \item[Threshold] $CDD \leq 0.10$ required
  \item[Drift Link] \driftid{TTS-DRIFT-005}
\end{description}

\subsubsection{TTS-M12: Taxonomy Coverage Drift (TCD)}

\begin{equation}
\boxed{
TCD = |LRC_{\text{baseline}} - LRC_{\text{current}}|
}
\end{equation}

\begin{description}
  \item[Variables] $LRC$ = Level Ratio Coverage from Chapter~\ref{chap:evaluation}
  \item[Direction] Lower is better
  \item[Threshold] $TCD \leq 0.05$ required (5\% coverage change tolerance)
  \item[Drift Link] \driftid{TTS-DRIFT-002}
\end{description}

\subsection{Statistical Foundation Metrics (TTS-M13 to TTS-M15)}

These metrics provide rigorous statistical foundations for drift detection,
enabling hypothesis testing and multi-metric validation.

\subsubsection{TTS-M13: Wasserstein Distance (W1D)}

The \textbf{Wasserstein distance} (Earth Mover's Distance) quantifies the
minimum ``work'' needed to transform the synthetic distribution $Q$ into
the real distribution $P$:

\begin{equation}
\boxed{
W_1(P, Q) = \int_{-\infty}^{\infty} |F_P(x) - F_Q(x)| \, dx
}
\end{equation}

\begin{description}
  \item[Variables] $F_P$, $F_Q$ = cumulative distribution functions of real and
        synthetic data respectively
  \item[Direction] Lower is better
  \item[Threshold] $W_{1,\text{norm}} \leq 0.10$ required (normalized to [0,1])
  \item[Implementation] \texttt{scipy.stats.wasserstein\_distance()}
  \item[Use Cases] Continuous variables, complexity scores, numeric columns
  \item[Drift Link] All drift types (cross-cutting)
\end{description}

\textbf{Normalization formula:}
\begin{equation}
W_{1,\text{norm}} = \frac{W_1(P, Q)}{\max(|x_{\max} - x_{\min}|, \epsilon)}
\end{equation}

\subsubsection{TTS-M14: Population Stability Index (PSI)}

The \textbf{PSI} measures how a distribution has shifted across predefined
bins, widely used in financial risk modeling for scorecard stability:

\begin{equation}
\boxed{
PSI = \sum_{i=1}^{k} (P_i - Q_i) \cdot \ln\left(\frac{P_i}{Q_i}\right)
}
\end{equation}

\begin{description}
  \item[Variables] $P_i$ = proportion in bin $i$ for real data;
        $Q_i$ = proportion in bin $i$ for synthetic data;
        $k$ = number of bins (typically 10)
  \item[Direction] Lower is better
  \item[Threshold] Industry standard:
    \begin{itemize}
      \item $PSI < 0.10$ --- No significant drift
      \item $0.10 \leq PSI < 0.25$ --- Moderate drift (investigate)
      \item $PSI \geq 0.25$ --- Significant drift (action required)
    \end{itemize}
  \item[Implementation] Custom implementation with Laplace smoothing
  \item[Use Cases] Categorical columns, discretized continuous variables, enums
  \item[Drift Link] \driftid{TTS-DRIFT-001}, \driftid{TTS-DRIFT-002}
\end{description}

\textbf{Laplace smoothing to avoid division by zero:}
\begin{equation}
P_i' = \frac{P_i \cdot n + \alpha}{n + k\alpha}, \quad \alpha = 0.001
\end{equation}

\subsubsection{TTS-M15: Kolmogorov-Smirnov Statistic (KSS)}

The \textbf{KS test} is a non-parametric test measuring the maximum distance
between two empirical CDFs:

\begin{equation}
\boxed{
D_{KS} = \sup_x |F_P(x) - F_Q(x)|
}
\end{equation}

\begin{description}
  \item[Variables] $F_P$, $F_Q$ = empirical CDFs; $\sup$ = supremum (maximum)
  \item[Direction] Lower $D_{KS}$ is better; higher p-value is better
  \item[Threshold] $p\text{-value} > 0.05$ required (no significant drift at 95\% confidence)
  \item[Implementation] \texttt{scipy.stats.ks\_2samp()}
  \item[Use Cases] Continuous variables, location and shape differences
  \item[Drift Link] \driftid{TTS-DRIFT-005}, \driftid{TTS-DRIFT-006}
\end{description}

\textbf{Decision rule:}
\begin{equation}
\text{Drift detected} \iff p\text{-value} \leq 0.05 \text{ (reject } H_0: P = Q\text{)}
\end{equation}

\subsection{Statistical Metric Comparison}

\begin{table}[htbp]
\centering
\footnotesize
\caption{Statistical Metrics Comparison}
\label{tab:stat-metrics-comparison}
\begin{tabularx}{\textwidth}{@{}l l l X l@{}}
\toprule
\textbf{Metric} & \textbf{Applicability} & \textbf{Sensitivity} & \textbf{Theoretical Basis} & \textbf{Speed} \\
\midrule
W1D & Continuous, ordinal & Proportional to shifts & Optimal transport theory & Medium \\
PSI & Categorical, binned & Practical thresholds & KL-divergence variant & Fast \\
KSS & Continuous & Location \& shape & Non-parametric CDF test & Fast \\
JSD (M07) & Both & Center-sensitive & Symmetric KL-divergence & Fast \\
\bottomrule
\end{tabularx}
\end{table}

\section{Composite TTS-EVAL Score}

The \textbf{TTS Evaluation Score} combines all drift metrics into a single
composite score, following the pattern established in singapore-ai-safety:

\begin{equation}
\boxed{
\begin{aligned}
TTS\_EVAL = 100 \times \Big(
  &0.15 \cdot SCR +
  0.10 \cdot (1 - SSS_{\text{norm}}) +
  0.10 \cdot FKCR \\
  &+ 0.15 \cdot SAS +
  0.10 \cdot IPR +
  0.15 \cdot ESR \\
  &+ 0.10 \cdot RFS +
  0.05 \cdot ARR
\Big) \\
- 100 \times \Big(
  &0.05 \cdot CTMR +
  0.03 \cdot TDS +
  0.02 \cdot CDD
\Big)
\end{aligned}
}
\label{eq:tts-eval}
\end{equation}

\noindent where $SSS_{\text{norm}} = \min(SSS / 14.0, 1.0)$ normalizes staleness
to a 0--1 range (14 days = maximum acceptable staleness).

\subsection{TTS-EVAL Interpretation}

\begin{table}[htbp]
\centering
\caption{TTS-EVAL Score Interpretation}
\begin{tabular}{l l l}
\toprule
\textbf{Score Band} & \textbf{Status} & \textbf{Action} \\
\midrule
$TTS\_EVAL \geq 85$ & \textcolor{green}{\textbf{GREEN}} & Production ready, no action \\
$70 \leq TTS\_EVAL < 85$ & \textcolor{orange}{\textbf{AMBER}} & Investigate, schedule remediation \\
$55 \leq TTS\_EVAL < 70$ & \textcolor{orange}{\textbf{AMBER}} & Urgent remediation required \\
$TTS\_EVAL < 55$ & \textcolor{red}{\textbf{RED}} & Kill-switch trigger, halt production \\
\bottomrule
\end{tabular}
\end{table}

\section{Statistical Evaluation Rules}

Following the singapore-ai-safety specification, all metrics use the
\textbf{95\% Wilson Lower Confidence Bound} for ``higher is better'' metrics:

\begin{equation}
\boxed{
LCB_{95} = \frac{\hat{p} + \frac{z^2}{2n} - z\sqrt{\frac{\hat{p}(1-\hat{p})}{n} + \frac{z^2}{4n^2}}}{1 + \frac{z^2}{n}}
\quad \text{where } z = 1.96
}
\end{equation}

For ``lower is better'' metrics (CTMR, TDS, CDD, SSS), use the
\textbf{95\% Wilson Upper Confidence Bound}.

\textbf{Production gate rule:} A metric passes only if the confidence-bounded
value meets the threshold.

\section{Minimum Sample Sizes}

\begin{table}[htbp]
\centering
\caption{Minimum Sample Sizes for TTS Drift Detection}
\label{tab:tts-sample-sizes}
\begin{tabular}{l l l l l}
\toprule
\textbf{Evaluation Context} & \textbf{Min Queries} & \textbf{Repetitions} & \textbf{Total} & \textbf{Frequency} \\
\midrule
Training Data Generation & 500 & 1 & 500 & Per batch \\
CI/CD Gate & 100 & 3 & 300 & Per PR \\
Production Monitoring & 1000 & 1 & 1000 & Daily \\
User Population Sampling & See Table~\ref{tab:user-sampling} & -- & -- & Rolling \\
\bottomrule
\end{tabular}
\end{table}

\section{User Population Sampling Strategy}

\begin{table}[htbp]
\centering
\caption{Stratified Sampling by User Population}
\label{tab:user-sampling}
\begin{tabularx}{\textwidth}{@{}l l l l X@{}}
\toprule
\textbf{Segment} & \textbf{Definition} & \textbf{Sample Rate} & \textbf{Frequency} & \textbf{Focus Metrics} \\
\midrule
Power Users & $>100$ queries/day & 100\% & Real-time & All TTS-M01--M12 \\
Regular Users & 10--100/day & 20\% & Hourly & TTS-M05, M06, M09, M10 \\
Occasional Users & $<10$/day & 5\% & Daily & TTS-M09, M10 \\
New Users & First 30 days & 50\% & Real-time & TTS-M06, M08 (ambiguity) \\
\bottomrule
\end{tabularx}
\end{table}

\begin{lstlisting}[language=jsonx,caption={User population sampler (simula\_tts\_drift\_sampler.py)}]
@dataclass
class SamplingConfig:
    """Configuration for stratified user sampling."""
    
    power_user_threshold: int = 100  # queries/day
    regular_user_min: int = 10
    new_user_window_days: int = 30
    
    sampling_rates: dict[str, float] = field(default_factory=lambda: {
        "power_user": 1.0,
        "regular_user": 0.2,
        "occasional_user": 0.05,
        "new_user": 0.5,
    })

class UserPopulationSampler:
    """Stratified sampling for drift detection across user populations."""
    
    def __init__(self, config: SamplingConfig):
        self.config = config
        self._user_query_counts: dict[str, int] = {}
        self._user_first_seen: dict[str, datetime] = {}
    
    def classify_user(self, user_id: str) -> str:
        """Classify user into population segment."""
        first_seen = self._user_first_seen.get(user_id)
        if first_seen and (datetime.now() - first_seen).days < self.config.new_user_window_days:
            return "new_user"
        
        daily_count = self._user_query_counts.get(user_id, 0)
        if daily_count > self.config.power_user_threshold:
            return "power_user"
        elif daily_count >= self.config.regular_user_min:
            return "regular_user"
        return "occasional_user"
    
    def should_sample(self, user_id: str) -> bool:
        """Determine if this query should be sampled for drift analysis."""
        segment = self.classify_user(user_id)
        rate = self.config.sampling_rates[segment]
        return random.random() < rate
\end{lstlisting}

\section{Training-Time Drift Evaluation}

During Simula training data generation, drift metrics are computed to ensure
generated examples remain aligned with the current HANA schema:

\begin{lstlisting}[language=jsonx,caption={Training-time drift evaluator (simula\_tts\_drift\_evaluator.py)}]
@dataclass
class TTSMetricReport:
    """Complete drift metric report."""
    
    # Schema drift metrics
    scr: float  # Schema Coverage Rate
    sss: float  # Schema Staleness Score
    ctmr: float  # Column Type Mismatch Rate
    fkcr: float  # Foreign Key Consistency Rate
    
    # Semantic drift metrics
    sas: float  # Semantic Alignment Score
    ipr: float  # Intent Preservation Rate
    tds: float  # Terminology Drift Score
    arr: float  # Ambiguity Resolution Rate
    
    # Generation quality metrics
    esr: float  # Execution Success Rate
    rfs: float  # Result Fidelity Score
    cdd: float  # Complexity Distribution Drift
    tcd: float  # Taxonomy Coverage Drift
    
    # Composite score
    tts_eval: float
    
    # Metadata
    evaluated_at: str
    sample_size: int
    confidence_level: float = 0.95

class TTSTrainingEvaluator:
    """Evaluate drift metrics during training data generation."""
    
    def __init__(
        self,
        schema_registry: SchemaRegistry,
        embedding_model: str = "sentence-transformers/all-MiniLM-L6-v2",
    ):
        self.schema_registry = schema_registry
        self.embedding_model = embedding_model
    
    async def evaluate_batch(
        self,
        examples: list[TrainingExample],
        baseline: Optional[TTSMetricReport] = None,
    ) -> TTSMetricReport:
        """
        Evaluate drift metrics for a batch of training examples.
        
        Args:
            examples: Training examples to evaluate
            baseline: Optional baseline for drift comparison
            
        Returns:
            Complete drift metric report
        """
        # Schema drift metrics
        scr = self._compute_schema_coverage(examples)
        sss = self._compute_schema_staleness()
        ctmr = self._compute_type_mismatch_rate(examples)
        fkcr = self._compute_fk_consistency(examples)
        
        # Semantic drift metrics
        sas = await self._compute_semantic_alignment(examples)
        ipr = await self._compute_intent_preservation(examples)
        tds = self._compute_terminology_drift(examples)
        arr = await self._compute_ambiguity_resolution(examples)
        
        # Generation quality metrics
        esr = self._compute_execution_success(examples)
        rfs = await self._compute_result_fidelity(examples)
        
        # Distribution drift (requires baseline)
        cdd = 0.0
        tcd = 0.0
        if baseline:
            cdd = self._compute_complexity_drift(examples, baseline)
            tcd = self._compute_taxonomy_drift(examples, baseline)
        
        # Composite score
        tts_eval = self._compute_tts_eval(
            scr, sss, ctmr, fkcr, sas, ipr, tds, arr, esr, rfs, cdd
        )
        
        return TTSMetricReport(
            scr=scr, sss=sss, ctmr=ctmr, fkcr=fkcr,
            sas=sas, ipr=ipr, tds=tds, arr=arr,
            esr=esr, rfs=rfs, cdd=cdd, tcd=tcd,
            tts_eval=tts_eval,
            evaluated_at=datetime.now().isoformat(),
            sample_size=len(examples),
        )
    
    def _compute_schema_coverage(self, examples: list[TrainingExample]) -> float:
        """Compute Schema Coverage Rate (SCR)."""
        schema_columns = set()
        for table in self.schema_registry.get_all_tables():
            for column in table.columns:
                schema_columns.add(f"{table.table_name}.{column.name}")
        
        referenced_columns = set()
        for example in examples:
            # Parse SQL to extract column references
            columns = self._extract_columns_from_sql(example.sql)
            referenced_columns.update(columns)
        
        if not schema_columns:
            return 1.0
        return len(referenced_columns & schema_columns) / len(schema_columns)
    
    def _compute_tts_eval(
        self,
        scr: float, sss: float, ctmr: float, fkcr: float,
        sas: float, ipr: float, tds: float, arr: float,
        esr: float, rfs: float, cdd: float,
    ) -> float:
        """Compute composite TTS-EVAL score."""
        # Normalize SSS to 0-1 range (14 days max)
        sss_norm = min(sss / 14.0, 1.0)
        
        # Positive contributions
        positive = (
            0.15 * scr +
            0.10 * (1 - sss_norm) +
            0.10 * fkcr +
            0.15 * sas +
            0.10 * ipr +
            0.15 * esr +
            0.10 * rfs +
            0.05 * arr
        )
        
        # Negative contributions (penalties)
        negative = (
            0.05 * ctmr +
            0.03 * tds +
            0.02 * cdd
        )
        
        return 100 * positive - 100 * negative
\end{lstlisting}

\section{Production Monitoring}

\begin{lstlisting}[language=jsonx,caption={Production drift monitor (simula\_tts\_drift\_monitor.py)}]
@dataclass
class DriftAlert:
    """Alert for detected drift."""
    
    alert_id: str
    severity: str  # "LOW", "MEDIUM", "HIGH", "CRITICAL"
    drift_type: str  # TTS-DRIFT-001 to TTS-DRIFT-006
    metric_code: str  # TTS-M01 to TTS-M12
    current_value: float
    threshold: float
    baseline_value: float
    delta: float
    timestamp: str
    user_segment: Optional[str] = None
    sample_queries: list[str] = field(default_factory=list)

class TTSProductionMonitor:
    """Monitor drift in production Text-to-SQL pipeline."""
    
    def __init__(
        self,
        baseline: TTSMetricReport,
        drift_threshold: float = 0.10,
        alert_callback: Optional[Callable[[DriftAlert], None]] = None,
    ):
        self.baseline = baseline
        self.drift_threshold = drift_threshold
        self.alert_callback = alert_callback
        self._rolling_metrics: dict[str, deque] = defaultdict(lambda: deque(maxlen=1000))
    
    async def monitor_query(
        self,
        prompt: str,
        generated_sql: str,
        execution_result: Optional[Any],
        user_id: str,
    ) -> Optional[DriftAlert]:
        """
        Monitor a single query for drift signals.
        
        Args:
            prompt: Natural language prompt
            generated_sql: Generated SQL query
            execution_result: Result of SQL execution (None if failed)
            user_id: User identifier for segmentation
            
        Returns:
            DriftAlert if drift detected, None otherwise
        """
        # Track execution success
        execution_success = execution_result is not None
        self._rolling_metrics["esr"].append(1.0 if execution_success else 0.0)
        
        # Compute rolling ESR
        current_esr = sum(self._rolling_metrics["esr"]) / len(self._rolling_metrics["esr"])
        esr_delta = abs(current_esr - self.baseline.esr)
        
        # Check for significant drift
        if esr_delta > self.drift_threshold:
            alert = DriftAlert(
                alert_id=str(uuid.uuid4()),
                severity="CRITICAL" if esr_delta > 0.2 else "HIGH",
                drift_type="TTS-DRIFT-006",
                metric_code="TTS-M09",
                current_value=current_esr,
                threshold=0.95,
                baseline_value=self.baseline.esr,
                delta=esr_delta,
                timestamp=datetime.now().isoformat(),
                user_segment=self._classify_user(user_id),
            )
            
            if self.alert_callback:
                self.alert_callback(alert)
            
            return alert
        
        return None
    
    def get_rolling_report(self) -> TTSMetricReport:
        """Get drift report based on rolling window of queries."""
        # Aggregate rolling metrics
        esr = sum(self._rolling_metrics["esr"]) / max(len(self._rolling_metrics["esr"]), 1)
        
        # Return partial report with available metrics
        return TTSMetricReport(
            scr=self.baseline.scr,  # Use baseline for schema metrics
            sss=self.baseline.sss,
            ctmr=self.baseline.ctmr,
            fkcr=self.baseline.fkcr,
            sas=self.baseline.sas,
            ipr=self.baseline.ipr,
            tds=self.baseline.tds,
            arr=self.baseline.arr,
            esr=esr,
            rfs=self.baseline.rfs,
            cdd=0.0,  # Requires full evaluation
            tcd=0.0,
            tts_eval=self._estimate_tts_eval(esr),
            evaluated_at=datetime.now().isoformat(),
            sample_size=len(self._rolling_metrics["esr"]),
        )
\end{lstlisting}

\section{CI/CD Integration}

\subsection{Training Data Generation Gate}

The drift evaluator integrates with Simula's generation pipeline as a quality gate:

\begin{lstlisting}[language=bash,caption={CI/CD gate for TTS drift}]
# In GitHub Actions workflow
- name: Run TTS Drift Evaluation
  run: |
    python -m src.training.pipeline.simula_tts_drift_evaluator \
      --examples data/hana_generated/*.jsonl \
      --schema-registry docs/schema/simula/hana-schema-entry.schema.json \
      --baseline data/baselines/tts_baseline.json \
      --output reports/tts_drift_report.json
    
- name: Check TTS-EVAL Score
  run: |
    TTS_EVAL=$(jq '.tts_eval' reports/tts_drift_report.json)
    if (( $(echo "$TTS_EVAL < 70" | bc -l) )); then
      echo "TTS-EVAL score $TTS_EVAL below threshold 70"
      exit 1
    fi
\end{lstlisting}

\subsection{Drift Alert Workflow}

\begin{table}[htbp]
\centering
\caption{Drift Alert Response Actions}
\begin{tabularx}{\textwidth}{@{}l l X l@{}}
\toprule
\textbf{Severity} & \textbf{TTS-EVAL} & \textbf{Action} & \textbf{Timeline} \\
\midrule
CRITICAL & $< 55$ & Kill-switch, halt production, page on-call & Immediate \\
HIGH & $55$--$70$ & Block deployments, create P1 ticket & 4 hours \\
MEDIUM & $70$--$85$ & Create P2 ticket, schedule remediation & 1 week \\
LOW & $\geq 85$ & Log for tracking, no action required & -- \\
\bottomrule
\end{tabularx}
\end{table}

\section{HANA Cloud Telemetry Tables}

Drift metrics are persisted to HANA Cloud for historical analysis:

\begin{lstlisting}[language=sql,caption={TTS Drift telemetry tables}]
-- See docs/hana-cloud/TTS_DRIFT_TABLES_DDL.sql for complete DDL

-- Drift metrics history
CREATE COLUMN TABLE "FINSIGHT_GOV"."TTS_DRIFT_METRICS" (
    "METRIC_ID" NVARCHAR(64) NOT NULL,
    "TIMESTAMP" TIMESTAMP NOT NULL,
    "METRIC_CODE" NVARCHAR(20) NOT NULL,  -- TTS-M01 to TTS-M12
    "VALUE" DECIMAL(10,6),
    "THRESHOLD" DECIMAL(10,6),
    "PASSED" BOOLEAN,
    "USER_SEGMENT" NVARCHAR(50),
    "BATCH_ID" NVARCHAR(64),
    PRIMARY KEY ("METRIC_ID")
);

-- Schema snapshots for staleness tracking
CREATE COLUMN TABLE "FINSIGHT_GOV"."TTS_SCHEMA_SNAPSHOTS" (
    "SNAPSHOT_ID" NVARCHAR(64) NOT NULL,
    "SCHEMA_NAME" NVARCHAR(100),
    "TABLE_NAME" NVARCHAR(100),
    "COLUMN_HASH" NVARCHAR(64),
    "SNAPSHOT_AT" TIMESTAMP,
    "TRAINING_SYNC_AT" TIMESTAMP,
    PRIMARY KEY ("SNAPSHOT_ID")
);
\end{lstlisting}

\section{Stochastic Surprise Gate (SSG) Framework}
\label{sec:ssg-framework}

The SSG framework provides a unified, statistically rigorous approach to drift detection
across all four drift areas. It treats each drift-sensitive artifact as a \textbf{time series}
and uses a calibrated surprise score (Z-score) based on a Hull-White mean-reverting
stochastic differential equation.

\subsection{Mathematical Foundation}

The SSG models the drift-sensitive observable $x_t$ as evolving according to:

\begin{equation}
\boxed{
dx_t = a(\theta(t) - x_t)dt + \sigma dW_t
}
\label{eq:hull-white-sde}
\end{equation}

\begin{description}
  \item[Variables] $a$ = mean reversion speed; $\theta(t)$ = time-dependent long-run mean;
        $\sigma$ = volatility; $W_t$ = Wiener process
  \item[Discretization] Euler-Maruyama: $\hat{x}_{t+1} = x_t + a(\theta(t) - x_t)\Delta t$
\end{description}

\subsection{Surprise Score Calculation}

\subsubsection{Step 1: Compute Residual (Innovation)}

After observing $x_t$, the residual is:

\begin{equation}
\epsilon_t = x_t - \hat{x}_t
\end{equation}

\subsubsection{Step 2: Update Volatility Estimate}

Use exponentially weighted moving variance:

\begin{equation}
\boxed{
\sigma_t^2 = \lambda \sigma_{t-1}^2 + (1-\lambda)\epsilon_t^2
}
\end{equation}

where $\lambda \in [0.95, 0.99]$ is the forgetting factor.

\subsubsection{Step 3: Compute Z-Score (Surprise)}

\begin{equation}
\boxed{
Z_t = \frac{|\epsilon_t|}{\max(\sigma_t, \sigma_{\min})}
}
\end{equation}

where $\sigma_{\min} = 10^{-9}$ prevents division by zero.

\subsubsection{Step 4: Apply Surprise Gate}

\begin{equation}
\text{Drift detected} \iff Z_t > Z_{\text{threshold}}
\end{equation}

Default $Z_{\text{threshold}} = 2.0$ (corresponds to $\approx 95\%$ confidence).

\subsubsection{Step 5: Convert to p-Value}

For statistical reporting, convert Z-score to two-tailed p-value:

\begin{equation}
\boxed{
p = \text{erfc}\left(\frac{Z_t}{\sqrt{2}}\right)
}
\end{equation}

\subsection{SSG Surprise Metrics (TTS-M16 to TTS-M19)}

\begin{table}[htbp]
\centering
\footnotesize
\caption{Stochastic Surprise Gate Metrics}
\label{tab:ssg-metrics}
\begin{tabularx}{\textwidth}{@{}l l X l@{}}
\toprule
\textbf{ID} & \textbf{Name} & \textbf{Observable} & \textbf{Application} \\
\midrule
\metricid{TTS-M16} & Codebase Surprise Score (CSS) & Architecture constraint violations & Agent code drift \\
\metricid{TTS-M17} & Schema Surprise Score (SSS-HW) & Schema dissimilarity count & HANA schema drift \\
\metricid{TTS-M18} & Prompt Surprise Score (PSS) & Embedding distance from centroid & User/agent prompts \\
\metricid{TTS-M19} & Token Surprise Score (TSS) & $-\log_2 P(\text{token}) / H$ & TOON token stream \\
\bottomrule
\end{tabularx}
\end{table}

\subsubsection{TTS-M16: Codebase Surprise Score (CSS)}

For AI coding agents, measure deviation from architectural specification:

\begin{equation}
\boxed{
Z_t^{\text{code}} = \frac{|e_t - \hat{e}_t|}{\sigma_t^{\text{code}}}
}
\end{equation}

\begin{description}
  \item[Observable $e_t$] Vector of: import violations, module boundary violations,
        file length violations, pattern mismatch counts
  \item[Predictive Model] Hull-White with $\theta = 0$ (ideal state has zero violations)
  \item[Threshold] $Z^{\text{code}} > 2.0$ triggers CI block
  \item[Mitigation] Refresh \texttt{AGENTS.md}, request human review
\end{description}

\subsubsection{TTS-M17: Schema Surprise Score (SSS-HW)}

For HANA schema drift against Simula generator expectations:

\begin{equation}
\boxed{
Z_t^{\text{schema}} = \frac{|d_t - \hat{d}_t|}{\sigma_t^{\text{schema}}}
}
\end{equation}

\begin{description}
  \item[Observable $d_t$] Dissimilarity vector: mismatched columns, type changes,
        missing indices, FK violations
  \item[Predictive Model] Hull-White with time-dependent $\theta(t) = \theta_0 + \beta t$
        to accommodate planned migrations
  \item[Threshold] $Z^{\text{schema}} > 2.0$ triggers generator freeze
  \item[Mitigation] Auto-adjust generator spec or emit alert for manual alignment
\end{description}

\subsubsection{TTS-M18: Prompt Surprise Score (PSS)}

For user/agent prompt drift from training distribution:

\begin{equation}
\boxed{
Z_t^{\text{prompt}} = \frac{\|\text{embedding}_t - \mathbb{E}[\text{embedding}_t]\|_2}{\sigma_t^{\text{prompt}}}
}
\end{equation}

\begin{description}
  \item[Observable] Euclidean distance between prompt embedding and training centroid
  \item[Predictive Model] Moving average of embeddings with mean-reverting dynamic
  \item[Threshold] $Z^{\text{prompt}} > 2.0$ indicates out-of-distribution prompt
  \item[Mitigation] Switch to fallback template, log for retraining
\end{description}

\subsubsection{TTS-M19: Token Surprise Score (TSS)}

For TOON token stream drift:

\begin{equation}
\boxed{
Z_t^{\text{toon}} = \frac{S_t - \mathbb{E}[S]}{\sigma_t^{\text{toon}}}
\quad \text{where} \quad
S_t = \frac{-\log_2 P(\text{token}_t)}{H}
}
\end{equation}

\begin{description}
  \item[Observable $S_t$] Normalized token surprise (entropy-scaled negative log probability)
  \item[Token Model] Adaptive n-gram with exponential forgetting
  \item[Threshold] $Z^{\text{toon}} > 2.0$ indicates unexpected structure
  \item[Mitigation] Trigger DSPy re-optimization, insert correction token
\end{description}

\subsection{Significance Classification}

\begin{table}[htbp]
\centering
\caption{Z-Score Significance Levels}
\label{tab:z-significance}
\begin{tabular}{l c c l}
\toprule
\textbf{Z-Score Range} & \textbf{p-Value} & \textbf{Classification} & \textbf{Action} \\
\midrule
$Z \leq 1.0$ & $\geq 0.317$ & Normal & No action \\
$1.0 < Z \leq 2.0$ & $0.046$--$0.317$ & Noteworthy & Log for monitoring \\
$2.0 < Z \leq 3.0$ & $0.003$--$0.046$ & Significant & Trigger mitigation \\
$Z > 3.0$ & $< 0.003$ & Extreme & Immediate escalation \\
\bottomrule
\end{tabular}
\end{table}

\subsection{Online Calibration}

The SSG uses recursive least squares (RLS) with exponential forgetting to adapt model parameters:

\begin{lstlisting}[language=jsonx,caption={Online calibrator for Hull-White parameters}]
@dataclass
class OnlineCalibrator:
    """Recursive calibrator for Hull-White model parameters."""
    
    forgetting_factor: float = 0.98  # lambda
    
    # Parameter estimates
    mean_reversion_estimate: float = 0.5
    long_run_mean_estimate: float = 0.0
    volatility_estimate: float = 0.1
    
    # RLS state
    _P: float = 1.0  # Covariance estimate
    
    def update(self, x_t: float, x_hat: float, epsilon: float):
        """Update parameter estimates using new observation."""
        # Update covariance with forgetting
        self._P = (1 / self.forgetting_factor) * self._P
        
        # Compute Kalman gain
        K = self._P / (1 + self._P)
        
        # Update parameter estimates
        self.mean_reversion_estimate += K * epsilon * (x_t - self.long_run_mean_estimate)
        
        # Update volatility using EWMA variance
        self.volatility_estimate = math.sqrt(
            self.forgetting_factor * self.volatility_estimate**2 +
            (1 - self.forgetting_factor) * epsilon**2
        )
        
        # Update covariance
        self._P = (1 - K) * self._P
\end{lstlisting}

\subsection{Audit Trail Structure}

Every surprise gate decision produces an immutable audit record:

\begin{lstlisting}[language=jsonx,caption={Surprise audit trail schema}]
@dataclass
class SurpriseAuditTrail:
    """Immutable audit record for surprise gate decisions."""
    
    # Identifiers
    audit_id: str
    timestamp: str
    observable_type: str  # "code", "schema", "prompt", "toon"
    
    # Observation
    observed_value: float
    predicted_value: float
    residual: float
    
    # Surprise metrics
    z_score: float
    p_value: float
    significance: str  # "normal", "noteworthy", "significant", "extreme"
    trigger: bool
    
    # Model state at decision time
    hull_white_state: dict[str, float]  # a, theta, sigma
    calibrator_version: str
    
    # Action taken
    mitigation_action: Optional[str]
    
    # Governance
    compliant_with: list[str]  # e.g., ["MAS-TRH-2024", "ISO-42001"]
\end{lstlisting}

\section{Future Directions: TOON and DSPy Integration}
\label{sec:tts-drift-future}

\begin{futurebox}
\textbf{v1.3+ Scope}\\
The following integrations are planned for future releases and are explicitly
out of scope for v1.2 implementation.
\end{futurebox}

\subsection{Token-Oriented Object Notation (TOON)}

TOON provides structured metadata that travels with AI-generated content.
For Text-to-SQL drift detection, TOON metadata enables:

\begin{itemize}
  \item \textbf{Provenance tracking} --- Schema version, training data version, model version
  \item \textbf{Drift signals} --- SCR at generation time, SAS at inference time
  \item \textbf{Governance} --- Drift score, last check timestamp, alert status
\end{itemize}

\begin{lstlisting}[language=jsonx,caption={TOON drift metadata structure}]
@dataclass
class TOONDriftMetadata:
    """TOON-compliant drift metadata for Text-to-SQL."""
    
    # Provenance
    schema_version: str
    training_data_version: str
    model_version: str
    
    # Drift Signals
    scr_at_generation: float
    sas_at_inference: float
    complexity_band: str  # "easy", "medium", "hard"
    
    # Traceability
    taxonomy_path: list[str]
    schema_tables_used: list[str]
    
    # Governance
    drift_score: float
    last_drift_check: str
    drift_alert_status: str  # "green", "amber", "red"
\end{lstlisting}

\subsection{DSPy Integration}

DSPy (Declarative Self-improving Language Programs) enables:

\begin{itemize}
  \item \textbf{Programmatic prompt optimization} --- Tune prompts based on drift metrics
  \item \textbf{Modular composition} --- Separate schema extraction, intent, SQL generation
  \item \textbf{Self-correction} --- Retry with refined prompts when drift detected
\end{itemize}

\begin{lstlisting}[language=jsonx,caption={DSPy module with drift guard (future)}]
import dspy

class TextToSQLWithDriftGuard(dspy.Module):
    """DSPy module with integrated drift detection."""
    
    def __init__(self, schema_registry: SchemaRegistry):
        super().__init__()
        self.schema_extractor = dspy.ChainOfThought(
            "question -> relevant_tables, relevant_columns"
        )
        self.sql_generator = dspy.ChainOfThought(
            "question, schema_context -> sql_query"
        )
        self.drift_detector = dspy.Predict(
            "sql_query, schema_context -> drift_score, drift_reasons"
        )
        self.schema_registry = schema_registry
    
    def forward(self, question: str):
        # Step 1: Schema extraction
        schema_pred = self.schema_extractor(question=question)
        
        # Step 2: SQL generation
        sql_pred = self.sql_generator(
            question=question,
            schema_context=self._get_context(schema_pred.relevant_tables)
        )
        
        # Step 3: Drift detection
        drift_pred = self.drift_detector(
            sql_query=sql_pred.sql_query,
            schema_context=self._get_context(schema_pred.relevant_tables)
        )
        
        # Step 4: Self-correction if drift detected
        if drift_pred.drift_score > 0.2:
            sql_pred = self._regenerate_with_guidance(
                question, drift_pred.drift_reasons
            )
        
        return dspy.Prediction(
            sql=sql_pred.sql_query,
            drift_score=drift_pred.drift_score,
        )
\end{lstlisting}

\section{Known Limitations \& Roadmap}

\subsection{Review Summary (v1.2)}

This specification has been reviewed against the following criteria.
\textbf{Overall Rating: 7.5/10} --- solid foundation with clear taxonomy and
mathematical rigor, but requires empirical validation before production deployment.

\begin{table}[htbp]
\centering
\caption{Specification Review Scores}
\label{tab:review-scores}
\begin{tabular}{l c l}
\toprule
\textbf{Criterion} & \textbf{Score} & \textbf{Assessment} \\
\midrule
Completeness & 7/10 & Missing latency/cost metrics \\
Mathematical Rigor & 9/10 & Strong statistical foundation \\
Practicality & 6/10 & Computational cost concerns \\
Innovation & 8/10 & SSG is novel \\
Governance \& Audit & 9/10 & Excellent audit trail design \\
Implementation Readiness & 5/10 & Needs empirical validation \\
Documentation Quality & 8/10 & Good cross-references \\
\bottomrule
\end{tabular}
\end{table}

\subsection{Known Limitations}

\begin{table}[htbp]
\centering
\footnotesize
\caption{Known Limitations and Mitigation Status}
\label{tab:limitations}
\begin{tabularx}{\textwidth}{@{}l X l l@{}}
\toprule
\textbf{Area} & \textbf{Issue} & \textbf{Severity} & \textbf{Status} \\
\midrule
Empirical validation & No results from real HANA data or agent logs & HIGH & v1.3 \\
Threshold justification & Thresholds not derived from error analysis & MEDIUM & v1.3 \\
Computational cost & No overhead analysis for real-time monitoring & MEDIUM & Addressed \\
User/agent classification & Heuristic rules are brittle & LOW & v1.3 \\
TOON + DSPy integration & Marked as future, conceptual only & MEDIUM & v1.3 \\
Ambiguity resolution (ARR) & Detection method not specified & HIGH & v1.3 \\
Calibration drift & No recovery from SSG calibration failure & MEDIUM & Addressed \\
Missing metrics & No latency drift or token cost drift & LOW & v2.0 \\
\bottomrule
\end{tabularx}
\end{table}

\subsection{Computational Optimizations}

To address the computational cost concerns, the following optimizations are
provided:

\subsubsection{Sliced Wasserstein Approximation}

For real-time monitoring, use the \textbf{sliced Wasserstein distance} which
projects distributions onto random 1-D subspaces:

\begin{equation}
\boxed{
SW_1(P, Q) = \frac{1}{L} \sum_{\ell=1}^{L} W_1(P_\ell, Q_\ell)
}
\end{equation}

where $P_\ell$ and $Q_\ell$ are 1-D projections onto random unit vectors.

\textbf{Complexity comparison:}
\begin{itemize}
  \item Full Wasserstein: $O(n^3 \log n)$ for exact, $O(n^2)$ for entropic
  \item Sliced Wasserstein: $O(L \cdot n \log n)$ where $L=50$ projections typical
\end{itemize}

\textbf{Implementation flag:}
\begin{lstlisting}[language=jsonx]
wasserstein_result = StatisticalDriftCalculator.compute_wasserstein(
    p_values, q_values,
    use_fast_approximation=True,  # Use sliced Wasserstein
    num_projections=50,
)
\end{lstlisting}

\subsubsection{Tiered Metric Strategy}

\begin{table}[htbp]
\centering
\caption{Tiered Metrics by Latency Budget}
\label{tab:tiered-metrics}
\begin{tabular}{l l l l}
\toprule
\textbf{Tier} & \textbf{Latency} & \textbf{Metrics} & \textbf{Use Case} \\
\midrule
Real-time & $<10$ms & ESR, rolling counts & Per-query monitoring \\
Fast & $<100$ms & PSI, KS, sliced Wasserstein & Batch (per minute) \\
Full & $<1$s & Full Wasserstein, SSG, SAS & Offline analysis \\
\bottomrule
\end{tabular}
\end{table}

\subsection{Calibration Health Checks}

To detect and recover from SSG calibration failure, the following checks are
implemented:

\begin{lstlisting}[language=jsonx,caption={Calibration health check}]
@dataclass
class CalibrationHealth:
    """Health status of SSG calibration."""
    
    volatility: float
    volatility_z_score: float
    is_healthy: bool
    recovery_action: Optional[str]

class HullWhiteSurpriseGate:
    
    SIGMA_MIN_THRESHOLD: float = 1e-6
    SIGMA_MAX_THRESHOLD: float = 100.0
    
    def check_calibration_health(self) -> CalibrationHealth:
        """Check if SSG calibration is healthy."""
        sigma = self.state.volatility
        
        # Check for implausibly low volatility (collapsed)
        if sigma < self.SIGMA_MIN_THRESHOLD:
            return CalibrationHealth(
                volatility=sigma,
                volatility_z_score=float("inf"),
                is_healthy=False,
                recovery_action="reset_calibrator_sigma_collapsed",
            )
        
        # Check for implausibly high volatility (exploded)
        if sigma > self.SIGMA_MAX_THRESHOLD:
            return CalibrationHealth(
                volatility=sigma,
                volatility_z_score=float("inf"),
                is_healthy=False,
                recovery_action="reset_calibrator_sigma_exploded",
            )
        
        return CalibrationHealth(
            volatility=sigma,
            volatility_z_score=0.0,
            is_healthy=True,
            recovery_action=None,
        )
    
    def auto_recover_calibration(self) -> bool:
        """Auto-recover from calibration failure."""
        health = self.check_calibration_health()
        if not health.is_healthy:
            self.calibrator.reset()
            self.state.volatility = self.initial_volatility
            return True
        return False
\end{lstlisting}

\subsection{Roadmap}

\begin{table}[htbp]
\centering
\footnotesize
\caption{Implementation Roadmap}
\label{tab:roadmap}
\begin{tabularx}{\textwidth}{@{}l l X l@{}}
\toprule
\textbf{Version} & \textbf{Target} & \textbf{Deliverables} & \textbf{Priority} \\
\midrule
v1.2.1 & 2026-Q2 & Sliced Wasserstein, calibration health checks & HIGH \\
v1.3 & 2026-Q3 & Spider validation, threshold bootstrap, ARR module & HIGH \\
v1.3 & 2026-Q3 & TOON/DSPy working implementation & MEDIUM \\
v1.3 & 2026-Q3 & User/agent classification improvement & LOW \\
v2.0 & 2026-Q4 & TTS-M20 (Latency), TTS-M21 (Token Cost) & MEDIUM \\
v2.0 & 2026-Q4 & Multivariate Mahalanobis surprise & LOW \\
\bottomrule
\end{tabularx}
\end{table}

\subsection{Pilot Deployment Strategy}

Before production deployment, conduct a pilot with the following parameters:

\begin{enumerate}
  \item \textbf{Traffic sample:} 1\% of production HANA queries
  \item \textbf{Metrics:} Lightweight only (PSI, KS, rolling ESR)
  \item \textbf{Heavy metrics:} Defer Wasserstein and full SSG to offline batch
  \item \textbf{Goal:} Collect empirical drift patterns for threshold tuning
  \item \textbf{Duration:} 4 weeks minimum
\end{enumerate}

\section{Detailed Implementation Specifications}

This section provides implementation-ready specifications for roadmap items.

\subsection{Sliced Wasserstein Distance Implementation}
\label{sec:sliced-wasserstein-impl}

\subsubsection{Algorithm}

\begin{algorithm}[H]
\caption{Sliced Wasserstein Distance}
\label{alg:sliced-wasserstein}
\begin{algorithmic}[1]
\Require Samples $P = \{p_1, \ldots, p_m\}$, $Q = \{q_1, \ldots, q_n\}$, projections $L$
\Ensure Sliced Wasserstein distance $SW_1(P, Q)$
\State $sw \gets 0$
\For{$\ell = 1$ to $L$}
    \State $\theta_\ell \gets \text{random\_unit\_vector}(d)$ \Comment{Sample direction}
    \State $P_\ell \gets \{\theta_\ell^\top p_i : p_i \in P\}$ \Comment{Project P}
    \State $Q_\ell \gets \{\theta_\ell^\top q_j : q_j \in Q\}$ \Comment{Project Q}
    \State $P_\ell^{\text{sorted}} \gets \text{sort}(P_\ell)$
    \State $Q_\ell^{\text{sorted}} \gets \text{sort}(Q_\ell)$
    \If{$m \neq n$} \Comment{Interpolate if sizes differ}
        \State $P_\ell^{\text{interp}}, Q_\ell^{\text{interp}} \gets \text{interpolate\_to\_common\_grid}(P_\ell, Q_\ell)$
        \State $w_\ell \gets \frac{1}{k}\sum_{i=1}^{k} |P_\ell^{\text{interp}}[i] - Q_\ell^{\text{interp}}[i]|$
    \Else
        \State $w_\ell \gets \frac{1}{m}\sum_{i=1}^{m} |P_\ell^{\text{sorted}}[i] - Q_\ell^{\text{sorted}}[i]|$
    \EndIf
    \State $sw \gets sw + w_\ell$
\EndFor
\State \Return $sw / L$
\end{algorithmic}
\end{algorithm}

\subsubsection{Python Implementation}

\begin{lstlisting}[language=jsonx,caption={Sliced Wasserstein implementation}]
import numpy as np
from typing import Optional

def compute_sliced_wasserstein(
    p_values: np.ndarray,
    q_values: np.ndarray,
    num_projections: int = 50,
    seed: Optional[int] = None,
) -> float:
    """
    Compute sliced Wasserstein distance.
    
    Args:
        p_values: Array of shape (m,) or (m, d) for d-dimensional data
        q_values: Array of shape (n,) or (n, d)
        num_projections: Number of random projections (L)
        seed: Random seed for reproducibility
        
    Returns:
        Sliced Wasserstein distance
    """
    rng = np.random.default_rng(seed)
    
    # Handle 1D case
    if p_values.ndim == 1:
        p_values = p_values.reshape(-1, 1)
    if q_values.ndim == 1:
        q_values = q_values.reshape(-1, 1)
    
    d = p_values.shape[1]
    m, n = len(p_values), len(q_values)
    
    # Common grid size for interpolation
    grid_size = max(m, n)
    
    total_distance = 0.0
    
    for _ in range(num_projections):
        # Sample random direction on unit sphere
        theta = rng.standard_normal(d)
        theta = theta / np.linalg.norm(theta)
        
        # Project samples
        p_proj = p_values @ theta
        q_proj = q_values @ theta
        
        # Sort projections
        p_sorted = np.sort(p_proj)
        q_sorted = np.sort(q_proj)
        
        # Interpolate to common grid if sizes differ
        if m != n:
            p_interp = np.interp(
                np.linspace(0, 1, grid_size),
                np.linspace(0, 1, m),
                p_sorted
            )
            q_interp = np.interp(
                np.linspace(0, 1, grid_size),
                np.linspace(0, 1, n),
                q_sorted
            )
            w_l = np.mean(np.abs(p_interp - q_interp))
        else:
            w_l = np.mean(np.abs(p_sorted - q_sorted))
        
        total_distance += w_l
    
    return total_distance / num_projections
\end{lstlisting}

\subsubsection{Parameters and Thresholds}

\begin{table}[htbp]
\centering
\caption{Sliced Wasserstein Parameters}
\begin{tabular}{l l l l}
\toprule
\textbf{Parameter} & \textbf{Default} & \textbf{Range} & \textbf{Rationale} \\
\midrule
\texttt{num\_projections} & 50 & 20--200 & 50 gives $<5\%$ variance \\
\texttt{seed} & None & Any int & Set for reproducibility in CI \\
\texttt{threshold} & 0.10 & 0.05--0.20 & Same as full W1D \\
\bottomrule
\end{tabular}
\end{table}

\subsection{Calibration Health Check Implementation}
\label{sec:calibration-health-impl}

\subsubsection{Health Check Algorithm}

\begin{algorithm}[H]
\caption{SSG Calibration Health Check}
\label{alg:calibration-health}
\begin{algorithmic}[1]
\Require SSG state $S = (a, \theta, \sigma_t)$, history window $H$
\Ensure Health status, recovery action
\State $\sigma_{\min} \gets 10^{-6}$, $\sigma_{\max} \gets 100.0$
\State $a_{\min} \gets 0.01$, $a_{\max} \gets 10.0$
\If{$\sigma_t < \sigma_{\min}$}
    \State \Return (\texttt{UNHEALTHY}, ``sigma\_collapsed'', \texttt{RESET})
\EndIf
\If{$\sigma_t > \sigma_{\max}$}
    \State \Return (\texttt{UNHEALTHY}, ``sigma\_exploded'', \texttt{RESET})
\EndIf
\If{$a < a_{\min}$ or $a > a_{\max}$}
    \State \Return (\texttt{UNHEALTHY}, ``mean\_reversion\_invalid'', \texttt{RESET})
\EndIf
\State $\sigma_{\text{history}} \gets [\sigma_{t-k} : k \in H]$ \Comment{Collect historical volatility}
\State $\mu_\sigma \gets \text{mean}(\sigma_{\text{history}})$
\State $\text{sd}_\sigma \gets \text{std}(\sigma_{\text{history}})$
\State $z_\sigma \gets |\sigma_t - \mu_\sigma| / \max(\text{sd}_\sigma, 10^{-9})$
\If{$z_\sigma > 3.0$}
    \State \Return (\texttt{WARNING}, ``sigma\_drifted'', \texttt{LOG\_ONLY})
\EndIf
\State \Return (\texttt{HEALTHY}, None, None)
\end{algorithmic}
\end{algorithm}

\subsubsection{Python Implementation}

\begin{lstlisting}[language=jsonx,caption={Calibration health check implementation}]
from dataclasses import dataclass
from enum import Enum
from typing import Optional, List
from collections import deque

class HealthStatus(Enum):
    HEALTHY = "healthy"
    WARNING = "warning"
    UNHEALTHY = "unhealthy"

class RecoveryAction(Enum):
    NONE = "none"
    LOG_ONLY = "log_only"
    RESET = "reset"
    ESCALATE = "escalate"

@dataclass
class CalibrationHealth:
    """Detailed health status of SSG calibration."""
    
    status: HealthStatus
    reason: Optional[str]
    action: RecoveryAction
    
    # Current state
    volatility: float
    mean_reversion: float
    
    # Diagnostic metrics
    volatility_z_score: float
    volatility_history_mean: float
    volatility_history_std: float
    observations_since_reset: int
    
    def to_audit_record(self) -> dict:
        """Convert to audit trail record."""
        return {
            "status": self.status.value,
            "reason": self.reason,
            "action": self.action.value,
            "volatility": self.volatility,
            "mean_reversion": self.mean_reversion,
            "volatility_z_score": self.volatility_z_score,
            "volatility_history_mean": self.volatility_history_mean,
            "volatility_history_std": self.volatility_history_std,
            "observations_since_reset": self.observations_since_reset,
        }

class CalibrationHealthChecker:
    """Check and maintain SSG calibration health."""
    
    # Thresholds
    SIGMA_MIN = 1e-6
    SIGMA_MAX = 100.0
    A_MIN = 0.01
    A_MAX = 10.0
    Z_WARNING_THRESHOLD = 3.0
    
    def __init__(self, history_window: int = 100):
        self.history_window = history_window
        self._sigma_history: deque = deque(maxlen=history_window)
        self._observations = 0
    
    def record_observation(self, sigma: float) -> None:
        """Record a volatility observation."""
        self._sigma_history.append(sigma)
        self._observations += 1
    
    def check_health(
        self,
        sigma: float,
        mean_reversion: float,
    ) -> CalibrationHealth:
        """
        Check calibration health.
        
        Args:
            sigma: Current volatility estimate
            mean_reversion: Current mean reversion speed
            
        Returns:
            CalibrationHealth with status and recommended action
        """
        # Check hard bounds on sigma
        if sigma < self.SIGMA_MIN:
            return CalibrationHealth(
                status=HealthStatus.UNHEALTHY,
                reason="sigma_collapsed",
                action=RecoveryAction.RESET,
                volatility=sigma,
                mean_reversion=mean_reversion,
                volatility_z_score=float("inf"),
                volatility_history_mean=0.0,
                volatility_history_std=0.0,
                observations_since_reset=self._observations,
            )
        
        if sigma > self.SIGMA_MAX:
            return CalibrationHealth(
                status=HealthStatus.UNHEALTHY,
                reason="sigma_exploded",
                action=RecoveryAction.RESET,
                volatility=sigma,
                mean_reversion=mean_reversion,
                volatility_z_score=float("inf"),
                volatility_history_mean=0.0,
                volatility_history_std=0.0,
                observations_since_reset=self._observations,
            )
        
        # Check mean reversion bounds
        if mean_reversion < self.A_MIN or mean_reversion > self.A_MAX:
            return CalibrationHealth(
                status=HealthStatus.UNHEALTHY,
                reason="mean_reversion_invalid",
                action=RecoveryAction.RESET,
                volatility=sigma,
                mean_reversion=mean_reversion,
                volatility_z_score=0.0,
                volatility_history_mean=0.0,
                volatility_history_std=0.0,
                observations_since_reset=self._observations,
            )
        
        # Check volatility stability (need history)
        if len(self._sigma_history) >= 10:
            import statistics
            mu = statistics.mean(self._sigma_history)
            sd = statistics.stdev(self._sigma_history) if len(self._sigma_history) > 1 else 1e-9
            z = abs(sigma - mu) / max(sd, 1e-9)
            
            if z > self.Z_WARNING_THRESHOLD:
                return CalibrationHealth(
                    status=HealthStatus.WARNING,
                    reason="sigma_drifted",
                    action=RecoveryAction.LOG_ONLY,
                    volatility=sigma,
                    mean_reversion=mean_reversion,
                    volatility_z_score=z,
                    volatility_history_mean=mu,
                    volatility_history_std=sd,
                    observations_since_reset=self._observations,
                )
            
            return CalibrationHealth(
                status=HealthStatus.HEALTHY,
                reason=None,
                action=RecoveryAction.NONE,
                volatility=sigma,
                mean_reversion=mean_reversion,
                volatility_z_score=z,
                volatility_history_mean=mu,
                volatility_history_std=sd,
                observations_since_reset=self._observations,
            )
        
        # Not enough history yet
        return CalibrationHealth(
            status=HealthStatus.HEALTHY,
            reason="insufficient_history",
            action=RecoveryAction.NONE,
            volatility=sigma,
            mean_reversion=mean_reversion,
            volatility_z_score=0.0,
            volatility_history_mean=0.0,
            volatility_history_std=0.0,
            observations_since_reset=self._observations,
        )
\end{lstlisting}

\subsection{Spider Dataset Validation Specification}
\label{sec:spider-validation-impl}

\subsubsection{Dataset Requirements}

\begin{table}[htbp]
\centering
\caption{Spider Dataset Validation Parameters}
\begin{tabular}{l l l}
\toprule
\textbf{Parameter} & \textbf{Value} & \textbf{Source} \\
\midrule
Dataset & Spider v1.0 & \url{https://yale-lily.github.io/spider} \\
Split & dev (1,034 examples) & Standard benchmark \\
Databases & 20 databases & Cross-domain coverage \\
Complexity levels & Easy, Medium, Hard, Extra & Per Spider taxonomy \\
\bottomrule
\end{tabular}
\end{table}

\subsubsection{Synthetic Drift Injection Protocol}

\begin{lstlisting}[language=jsonx,caption={Drift injection for validation}]
@dataclass
class DriftInjectionConfig:
    """Configuration for synthetic drift injection."""
    
    # Schema drift parameters
    column_drop_rate: float = 0.05  # 5% of columns dropped
    column_rename_rate: float = 0.10  # 10% renamed
    type_change_rate: float = 0.03  # 3% type changes
    
    # Prompt drift parameters
    vocabulary_shift_rate: float = 0.15  # 15% new terms
    complexity_shift: float = 0.20  # 20% harder queries
    
    # Temporal parameters
    drift_start_fraction: float = 0.30  # Drift starts at 30%
    drift_ramp_fraction: float = 0.20  # Full drift by 50%

class SpiderDriftInjector:
    """Inject synthetic drift into Spider dataset."""
    
    def __init__(self, config: DriftInjectionConfig):
        self.config = config
        self._original_schemas: dict[str, dict] = {}
    
    def inject_schema_drift(
        self,
        schema: dict,
        drift_intensity: float,
    ) -> tuple[dict, list[str]]:
        """
        Inject schema drift.
        
        Args:
            schema: Original database schema
            drift_intensity: 0.0 (no drift) to 1.0 (full drift)
            
        Returns:
            (modified_schema, list_of_changes)
        """
        changes = []
        modified = copy.deepcopy(schema)
        
        for table in modified["tables"]:
            for col in table["columns"]:
                r = random.random()
                
                # Column drop
                if r < self.config.column_drop_rate * drift_intensity:
                    col["_dropped"] = True
                    changes.append(f"DROP {table['name']}.{col['name']}")
                    continue
                
                # Column rename
                if r < (self.config.column_drop_rate + self.config.column_rename_rate) * drift_intensity:
                    old_name = col["name"]
                    col["name"] = f"{old_name}_v2"
                    changes.append(f"RENAME {table['name']}.{old_name} -> {col['name']}")
                    continue
                
                # Type change
                if r < (self.config.column_drop_rate + self.config.column_rename_rate + 
                        self.config.type_change_rate) * drift_intensity:
                    old_type = col["type"]
                    col["type"] = self._mutate_type(old_type)
                    changes.append(f"TYPE {table['name']}.{col['name']}: {old_type} -> {col['type']}")
        
        # Filter dropped columns
        for table in modified["tables"]:
            table["columns"] = [c for c in table["columns"] if not c.get("_dropped")]
        
        return modified, changes
    
    def inject_prompt_drift(
        self,
        prompt: str,
        drift_intensity: float,
    ) -> str:
        """Inject vocabulary drift into prompt."""
        words = prompt.split()
        for i, word in enumerate(words):
            if random.random() < self.config.vocabulary_shift_rate * drift_intensity:
                words[i] = self._synonym_or_noise(word)
        return " ".join(words)
\end{lstlisting}

\subsubsection{Validation Metrics Protocol}

\begin{enumerate}
  \item \textbf{Baseline Phase} (examples 1--300):
    \begin{itemize}
      \item No drift injected
      \item Compute baseline TTS-EVAL, store as $TTS\_EVAL_{\text{baseline}}$
      \item Record all 19 metric baselines
    \end{itemize}
  
  \item \textbf{Drift Ramp Phase} (examples 301--500):
    \begin{itemize}
      \item Linearly increase drift intensity from 0.0 to 1.0
      \item Record TTS-EVAL at each 50-example checkpoint
      \item Record first-detection example index for each metric
    \end{itemize}
  
  \item \textbf{Full Drift Phase} (examples 501--1034):
    \begin{itemize}
      \item Drift intensity = 1.0 (constant)
      \item Compute final TTS-EVAL, compare to baseline
      \item Record false positive rate (alerts during baseline)
      \item Record false negative rate (missed drift during ramp)
    \end{itemize}
\end{enumerate}

\subsubsection{Expected Outputs}

\begin{table}[htbp]
\centering
\caption{Spider Validation Expected Outputs}
\begin{tabular}{l l l}
\toprule
\textbf{Output} & \textbf{File} & \textbf{Contents} \\
\midrule
Baseline report & \texttt{spider\_baseline.json} & All 19 metrics at t=300 \\
Drift detection curve & \texttt{spider\_drift\_curve.csv} & TTS-EVAL vs example index \\
False positive rate & \texttt{spider\_fpr.json} & FPR per metric (should be $<5\%$) \\
False negative rate & \texttt{spider\_fnr.json} & FNR per metric (should be $<10\%$) \\
Latency report & \texttt{spider\_latency.json} & P50, P95, P99 per metric \\
\bottomrule
\end{tabular}
\end{table}

\subsection{Threshold Bootstrap Analysis Specification}
\label{sec:threshold-bootstrap-impl}

\subsubsection{Bootstrap Protocol}

\begin{algorithm}[H]
\caption{Threshold Bootstrap via Percentile Method}
\label{alg:threshold-bootstrap}
\begin{algorithmic}[1]
\Require Clean dataset $D$ (no drift), metric function $f$, target FPR $\alpha$
\Ensure Threshold $\tau$ such that $P(f(x) > \tau | \text{no drift}) \le \alpha$
\State $B \gets 1000$ \Comment{Bootstrap iterations}
\State $n \gets |D|$
\State $S \gets []$ \Comment{Bootstrap metric samples}
\For{$b = 1$ to $B$}
    \State $D^*_b \gets \text{sample\_with\_replacement}(D, n)$
    \State $m_b \gets f(D^*_b)$ \Comment{Compute metric on bootstrap sample}
    \State $S.\text{append}(m_b)$
\EndFor
\If{metric direction is ``higher is better''}
    \State $\tau \gets \text{percentile}(S, \alpha \times 100)$ \Comment{Lower tail}
\Else
    \State $\tau \gets \text{percentile}(S, (1 - \alpha) \times 100)$ \Comment{Upper tail}
\EndIf
\State \Return $\tau$
\end{algorithmic}
\end{algorithm}

\subsubsection{Python Implementation}

\begin{lstlisting}[language=jsonx,caption={Threshold bootstrap implementation}]
import numpy as np
from typing import Callable, List, Tuple
from dataclasses import dataclass

@dataclass
class BootstrapResult:
    """Result of bootstrap threshold analysis."""
    
    metric_name: str
    direction: str  # "higher_better" or "lower_better"
    
    # Bootstrap distribution
    bootstrap_mean: float
    bootstrap_std: float
    bootstrap_percentiles: dict[int, float]  # {5: 0.82, 50: 0.90, 95: 0.95}
    
    # Derived thresholds at different FPR levels
    threshold_fpr_01: float  # 1% FPR
    threshold_fpr_05: float  # 5% FPR
    threshold_fpr_10: float  # 10% FPR
    
    # Validation
    actual_fpr: float  # Measured FPR on holdout
    num_bootstrap_samples: int

def bootstrap_threshold(
    data: List[dict],
    metric_fn: Callable[[List[dict]], float],
    metric_name: str,
    direction: str,
    num_bootstrap: int = 1000,
    target_fpr: float = 0.05,
    seed: int = 42,
) -> BootstrapResult:
    """
    Compute threshold via bootstrap.
    
    Args:
        data: Clean dataset (no drift)
        metric_fn: Function to compute metric from data
        metric_name: Name of the metric
        direction: "higher_better" or "lower_better"
        num_bootstrap: Number of bootstrap iterations
        target_fpr: Target false positive rate
        seed: Random seed
        
    Returns:
        BootstrapResult with threshold and diagnostics
    """
    rng = np.random.default_rng(seed)
    n = len(data)
    
    # Bootstrap sampling
    bootstrap_metrics = []
    for _ in range(num_bootstrap):
        indices = rng.choice(n, size=n, replace=True)
        sample = [data[i] for i in indices]
        m = metric_fn(sample)
        bootstrap_metrics.append(m)
    
    bootstrap_metrics = np.array(bootstrap_metrics)
    
    # Compute percentiles
    percentiles = {
        1: np.percentile(bootstrap_metrics, 1),
        5: np.percentile(bootstrap_metrics, 5),
        10: np.percentile(bootstrap_metrics, 10),
        25: np.percentile(bootstrap_metrics, 25),
        50: np.percentile(bootstrap_metrics, 50),
        75: np.percentile(bootstrap_metrics, 75),
        90: np.percentile(bootstrap_metrics, 90),
        95: np.percentile(bootstrap_metrics, 95),
        99: np.percentile(bootstrap_metrics, 99),
    }
    
    # Derive thresholds based on direction
    if direction == "higher_better":
        # Threshold is lower bound (alert if below)
        threshold_01 = np.percentile(bootstrap_metrics, 1)
        threshold_05 = np.percentile(bootstrap_metrics, 5)
        threshold_10 = np.percentile(bootstrap_metrics, 10)
    else:
        # Threshold is upper bound (alert if above)
        threshold_01 = np.percentile(bootstrap_metrics, 99)
        threshold_05 = np.percentile(bootstrap_metrics, 95)
        threshold_10 = np.percentile(bootstrap_metrics, 90)
    
    return BootstrapResult(
        metric_name=metric_name,
        direction=direction,
        bootstrap_mean=float(np.mean(bootstrap_metrics)),
        bootstrap_std=float(np.std(bootstrap_metrics)),
        bootstrap_percentiles=percentiles,
        threshold_fpr_01=threshold_01,
        threshold_fpr_05=threshold_05,
        threshold_fpr_10=threshold_10,
        actual_fpr=0.0,  # To be filled after validation
        num_bootstrap_samples=num_bootstrap,
    )
\end{lstlisting}

\subsubsection{Threshold Documentation Template}

Each threshold must be documented with:

\begin{lstlisting}[language=jsonx,caption={Threshold documentation schema}]
{
  "metric_code": "TTS-M01",
  "metric_name": "Schema Coverage Rate",
  "current_threshold": 0.85,
  "threshold_derivation": {
    "method": "bootstrap_percentile",
    "dataset": "Spider v1.0 dev split (n=1034)",
    "bootstrap_iterations": 1000,
    "target_fpr": 0.05,
    "bootstrap_mean": 0.912,
    "bootstrap_std": 0.023,
    "derived_threshold": 0.867,
    "chosen_threshold": 0.85,
    "rationale": "Rounded down to 0.85 for operational margin"
  },
  "validation": {
    "holdout_dataset": "Spider train subset (n=500)",
    "measured_fpr": 0.042,
    "measured_fnr": 0.08,
    "validated_at": "2026-04-23"
  }
}
\end{lstlisting}

\subsection{Ambiguity Resolution Rate (ARR) Module Specification}
\label{sec:arr-impl}

\subsubsection{Ambiguity Classification}

\begin{table}[htbp]
\centering
\footnotesize
\caption{Ambiguity Types for Text-to-SQL}
\begin{tabularx}{\textwidth}{@{}l X l@{}}
\toprule
\textbf{Type} & \textbf{Example} & \textbf{Resolution Strategy} \\
\midrule
Lexical & ``Show sales'' (which sales table?) & Schema lookup, clarification \\
Scope & ``Show recent orders'' (what is recent?) & Default + flag \\
Aggregation & ``Total revenue'' (sum? count?) & Context inference \\
Temporal & ``Last month'' (relative to when?) & Timestamp context \\
Entity & ``Show customer info'' (which customer?) & Require filter \\
\bottomrule
\end{tabularx}
\end{table}

\subsubsection{Ambiguity Detector Architecture}

\begin{lstlisting}[language=jsonx,caption={ARR module architecture}]
from dataclasses import dataclass
from enum import Enum
from typing import List, Optional

class AmbiguityType(Enum):
    LEXICAL = "lexical"
    SCOPE = "scope"
    AGGREGATION = "aggregation"
    TEMPORAL = "temporal"
    ENTITY = "entity"
    NONE = "none"

@dataclass
class AmbiguityDetection:
    """Result of ambiguity detection."""
    
    is_ambiguous: bool
    ambiguity_type: AmbiguityType
    confidence: float
    ambiguous_span: Optional[str]
    suggested_clarifications: List[str]
    resolution_applied: Optional[str]

class AmbiguityDetector:
    """
    Detect and classify ambiguity in Text-to-SQL prompts.
    
    Uses a two-stage approach:
    1. Rule-based heuristics for common patterns
    2. Small classifier for edge cases
    """
    
    # Heuristic patterns
    LEXICAL_PATTERNS = [
        r"\b(show|get|display)\s+\w+\b",  # Vague verbs
        r"\b(info|data|stuff)\b",  # Vague nouns
    ]
    
    SCOPE_PATTERNS = [
        r"\b(recent|latest|new|old)\b",  # Relative time
        r"\b(some|few|many|all)\b",  # Vague quantities
    ]
    
    TEMPORAL_PATTERNS = [
        r"\b(last|this|next)\s+(week|month|year|quarter)\b",
        r"\b(yesterday|today|tomorrow)\b",
    ]
    
    def __init__(
        self,
        schema_terms: set[str],
        classifier_model: Optional[str] = "microsoft/deberta-v3-small",
        classifier_threshold: float = 0.7,
    ):
        """
        Initialize detector.
        
        Args:
            schema_terms: Set of valid schema terms (tables, columns)
            classifier_model: HuggingFace model for edge cases
            classifier_threshold: Confidence threshold for classifier
        """
        self.schema_terms = {t.lower() for t in schema_terms}
        self.classifier_model = classifier_model
        self.classifier_threshold = classifier_threshold
        self._classifier = None  # Lazy load
    
    def detect(self, prompt: str, schema_context: dict) -> AmbiguityDetection:
        """
        Detect ambiguity in a prompt.
        
        Args:
            prompt: Natural language query
            schema_context: Available schema information
            
        Returns:
            AmbiguityDetection result
        """
        import re
        
        prompt_lower = prompt.lower()
        
        # Check lexical ambiguity (unknown terms)
        words = set(re.findall(r'\b\w+\b', prompt_lower))
        unknown_terms = words - self.schema_terms - self._stopwords()
        
        if len(unknown_terms) > 3:  # Many unknown terms
            return AmbiguityDetection(
                is_ambiguous=True,
                ambiguity_type=AmbiguityType.LEXICAL,
                confidence=0.8,
                ambiguous_span=", ".join(list(unknown_terms)[:3]),
                suggested_clarifications=[
                    f"Did you mean: {self._suggest_term(t)}" for t in list(unknown_terms)[:2]
                ],
                resolution_applied=None,
            )
        
        # Check scope ambiguity
        for pattern in self.SCOPE_PATTERNS:
            match = re.search(pattern, prompt_lower)
            if match:
                return AmbiguityDetection(
                    is_ambiguous=True,
                    ambiguity_type=AmbiguityType.SCOPE,
                    confidence=0.7,
                    ambiguous_span=match.group(),
                    suggested_clarifications=[
                        "Please specify an exact count or range",
                        "Consider adding a numeric limit"
                    ],
                    resolution_applied="default_limit_100",
                )
        
        # Check temporal ambiguity
        for pattern in self.TEMPORAL_PATTERNS:
            match = re.search(pattern, prompt_lower)
            if match:
                return AmbiguityDetection(
                    is_ambiguous=True,
                    ambiguity_type=AmbiguityType.TEMPORAL,
                    confidence=0.9,
                    ambiguous_span=match.group(),
                    suggested_clarifications=[
                        "Please specify exact dates",
                        f"Interpreting '{match.group()}' relative to current date"
                    ],
                    resolution_applied="relative_to_current_timestamp",
                )
        
        # Fall back to classifier for edge cases
        if self._classifier is None:
            self._load_classifier()
        
        if self._classifier:
            result = self._classifier_predict(prompt)
            if result["is_ambiguous"] and result["confidence"] > self.classifier_threshold:
                return AmbiguityDetection(
                    is_ambiguous=True,
                    ambiguity_type=AmbiguityType(result["type"]),
                    confidence=result["confidence"],
                    ambiguous_span=None,
                    suggested_clarifications=["Consider providing more specific terms"],
                    resolution_applied=None,
                )
        
        # No ambiguity detected
        return AmbiguityDetection(
            is_ambiguous=False,
            ambiguity_type=AmbiguityType.NONE,
            confidence=0.95,
            ambiguous_span=None,
            suggested_clarifications=[],
            resolution_applied=None,
        )
\end{lstlisting}

\subsubsection{ARR Metric Computation}

\begin{lstlisting}[language=jsonx,caption={ARR metric computation}]
def compute_arr(
    prompts: List[str],
    schema_context: dict,
    ground_truth_resolutions: List[str],
    detector: AmbiguityDetector,
) -> float:
    """
    Compute Ambiguity Resolution Rate.
    
    Args:
        prompts: List of prompts to evaluate
        schema_context: Schema information
        ground_truth_resolutions: Expected SQL for each prompt
        detector: Ambiguity detector instance
        
    Returns:
        ARR score in [0, 1]
    """
    ambiguous_count = 0
    correctly_resolved = 0
    
    for prompt, expected_sql in zip(prompts, ground_truth_resolutions):
        detection = detector.detect(prompt, schema_context)
        
        if detection.is_ambiguous:
            ambiguous_count += 1
            
            # Check if resolution was correct
            if detection.resolution_applied:
                # Generate SQL with resolution and compare
                resolved_sql = generate_sql_with_resolution(
                    prompt, detection.resolution_applied
                )
                if sql_equivalent(resolved_sql, expected_sql):
                    correctly_resolved += 1
    
    if ambiguous_count == 0:
        return 1.0  # No ambiguous prompts = perfect score
    
    return correctly_resolved / ambiguous_count
\end{lstlisting}

\section{Implementation Checklist}

\begin{itemize}
  \item[\texttt{[x]}] Define 19 TTS drift metrics (TTS-M01 to TTS-M19)
  \item[\texttt{[x]}] Specify calculation formulas with statistical bounds
  \item[\texttt{[x]}] Define TTS-EVAL composite score
  \item[\texttt{[x]}] Specify user population sampling strategy
  \item[\texttt{[x]}] Implement \texttt{simula\_tts\_drift\_evaluator.py}
  \item[\texttt{[x]}] Implement \texttt{simula\_tts\_drift\_monitor.py}
  \item[\texttt{[x]}] Implement \texttt{simula\_tts\_drift\_sampler.py}
  \item[\texttt{[x]}] Create HANA Cloud DDL for drift telemetry
  \item[\texttt{[x]}] Create JSON schemas for drift artifacts
  \item[\texttt{[x]}] Update spec-code-mapping.yaml
  \item[\texttt{[x]}] Implement HullWhiteSurpriseGate (SSG)
  \item[\texttt{[x]}] Add calibration health checks
  \item[\texttt{[ ]}] Integrate with CI/CD pipeline
  \item[\texttt{[ ]}] Spider dataset validation
  \item[\texttt{[ ]}] Threshold bootstrap analysis
  \item[\texttt{[ ]}] Ambiguity resolution module (ARR)
\end{itemize}
